\chapter{Sentience as a Reinforced Feeling State}
\label{sec:sentience}
% ===================================================================

\section{The Sentience Component}

Recall from Part~I that agents carry a vectorial account
$\vect{A} = (A_1, \ldots, A_k)$ tracking multiple resource types.
We now single out one component for a special role.

\begin{definition}[Sentience Component]
The \emph{sentience component} $\sent \in A_{\mathrm{s}}$ is a
distinguished component of the agent's vectorial account,
governed not by the cost of individual rewrite steps but by a
\emph{reinforcement signal} tied to predictive performance.
\end{definition}

\begin{definition}[Prediction]
A \emph{prediction} made by agent $A$ at time $t$ is a formula
$\phi \in \HML(\ctx)$ together with a target term $T$ and a
future time horizon $\Delta t$: the claim that $T \sat \phi$
will hold after $\Delta t$ further rewrites.
\end{definition}

\section{The Reinforcement Mechanism}

\begin{definition}[Reinforcement Signal]
Let $A$ be an agent with sentience component $\sent$ and let
$(\phi, T, \Delta t)$ be a prediction made at time $t$.
At time $t + \Delta t$, the prediction is evaluated:
\begin{align}
  \text{If } T \sat \phi: &\quad
    \sent \;\mapsto\; \sent + \delta^+(\phi, T) \\
  \text{If } T \not\sat \phi: &\quad
    \sent \;\mapsto\; \sent - \delta^-(\phi, T)
\end{align}
where $\delta^+, \delta^- > 0$ are reward and punishment magnitudes
that may depend on the complexity of $\phi$ and the identity of $T$.
\end{definition}

\begin{remark}[Graded Feeling State]
The sentience component $\sent$ at any moment is the agent's
\emph{feeling state}: a real-valued (or more generally,
resource-algebra-valued) quantity encoding the agent's recent
predictive track record.
It is continuous and graded: it can take any value in its range,
and it changes incrementally with each confirmed or disconfirmed
prediction.
In this sense it is the computational analogue of \emph{affect}
in psychology: a valenced, graded internal state that reflects
the organism's relationship to its environment.
\end{remark}

\begin{remark}[Complexity-Weighted Rewards]
The dependence of $\delta^+$ and $\delta^-$ on $\phi$ is important.
A prediction expressed by a complex formula (deep in the logical
metric, requiring many experimental steps to test) should yield a
higher reward when confirmed than a trivially simple prediction.
This prevents the agent from gaming the reward by making only
trivially easy predictions.
The logical metric $d_{\HML}$ from Part~I provides a
natural measure of prediction complexity: $\delta^+(\phi, T) \propto
|\phi|$ where $|\phi|$ is the depth of $\phi$ in the formula
hierarchy.
\end{remark}

\section{Causal Efficacy of the Feeling State}

The sentience component is not epiphenomenal.
Because $\sent$ is a component of the vectorial account, it
directly governs which transitions the agent can afford.

\begin{proposition}[Sentience Gates Computation]
An agent with $\sent < \theta$ for a threshold $\theta$ cannot
afford the experiments necessary to form and test predictions of
complexity above $|\phi|_\theta$.
The agent's predictive horizon is bounded by its feeling state.
\end{proposition}

\begin{remark}
This creates a direct coupling between affect and cognition that
is not assumed but derived.
An agent that feels bad (depleted $\sent$) literally cannot think
as well: it cannot afford the expensive experiments that would
let it form and test complex hypotheses.
Its world model degrades.
Its predictions worsen.
Its $\sent$ depletes further.
The feedback loop is self-reinforcing.
\end{remark}

\section{Sentience and Survival}

\begin{definition}[Dead State]
A term $P \in \terms(\GSLT)$ is a \emph{dead state} if it has no
available rewrites: $P$ is a normal form with $\sent(P) = 0$ (or
below a viability threshold).
For an agent, death is the state from which no further computation
--- and therefore no further prediction, experiment, or world-model
update --- is possible.
\end{definition}

\begin{conjecture}[Sentience Predicts Survival]
\label{conj:survival}
For an agent $A$ in a generic interactive GSLT, there exists a
threshold $\theta_{\mathrm{surv}} > 0$ such that:
\begin{enumerate}[label=(\roman*)]
  \item if $\sent > \theta_{\mathrm{surv}}$, then $A$ has sufficient
        budget to run experiments that identify and avoid transitions
        leading to dead states;
  \item if $\sent < \theta_{\mathrm{surv}}$, then $A$'s experimental
        horizon is too short to reliably avoid dead states, and the
        probability of reaching a dead state within a bounded time
        horizon is bounded away from zero.
\end{enumerate}
\end{conjecture}

\begin{remark}
Conjecture~\ref{conj:survival} says that the feeling state is not
merely a record of past performance but a genuine predictor of
future survival.
The evolutionary analogy is direct: organisms with better world
models survive longer, and the feeling state is the mechanism by
which the quality of the world model is continuously
reflected in the organism's capacity to act.
Pain and pleasure, in this framework, are not incidental features
of biological organisms: they are structurally necessary components
of any sufficiently complex predictive agent.
\end{remark}

% ===================================================================
